% element: 10-s062:e04
\BeleskeID{10-s062}
U aktivnom režimu kolektorska struja može da se izrazi i preko emitorske:
\[I_{C1}=\frac{\beta_F}{\beta_F+1}I_{E1}.\]
Kolektor–emitor napon dobijamo kao $V_{CE1}=V_{C1}-V_{E1}$. Zamenom izraza za oba čvora i uslova $V_{CE1}>V_{CES}$ dolazimo do dve nejednakosti na slajdu. Smer nejednakosti zavisi od znaka faktora $1-R_C\beta_F/R_B$.

Za izlazni opseg $0\leq V_O\leq V_{CC}$ režim zavisi od odnosa baznog i kolektorskog otpornika. Ako je $R_C=0$, tranzistor uvek radi aktivno. Ako je $\beta_FR_C<R_B$, praktično ostaje u aktivnom režimu. Pri obrnutom odnosu, zavisno od konkretnih vrednosti, viši izlazni naponi mogu odgovarati aktivnom režimu, a niži zasićenju.
